%%% Local Variables:
%%% mode: latex
%%% TeX-master: "../doc"
%%% coding: utf-8
%%% End:
% !TEX TS-program = pdflatexmk
% !TEX encoding = UTF-8 Unicode
% !TEX root = ../doc.tex

% TODO: 4 or 5 dimensions of trust?
 
\begin{quote}
\textit{``Is it [the LLM-Model's output] not only correct, but actually, can I trust it?''}\\
\hspace*{\fill}--- Thomas Kuster, Partner at LEXR, GenAI ZH 2026
\end{quote}
 
\noindent Large Language Models (LLMs) such as ChatGPT, Claude, and Gemini have moved from
research curiosity to everyday tool with remarkable speed.
Since the public release of ChatGPT in November 2022, awareness of AI-powered
applications has spread widely: a 2025 survey by gfs.bern commissioned by the Swiss
Federal Chancellery found that 96\% of the Swiss population had heard of
LLM-based tools such as ChatGPT, Gemini, or Copilot \cite{bern_chance_2025}.
Of those familiar with these applications, 59\% reported using them at least
occasionally, and 59\% of users also employ them in a professional context
\cite{bern_chance_2025}.
These figures illustrate that LLMs are no longer a niche technology; instead they are embedded
in the daily routines of a broad and diverse population, the majority of whom are not
AI specialists.
 
Widespread adoption, however, does not imply well-calibrated trust.
LLMs are susceptible to two distinct failure modes (dishonesty, in which a model
contradicts knowledge it has acquired during training, and hallucination, in which it
fabricates information entirely \cite{wu_usable_2025}), yet their outputs are delivered
with the same confident fluency regardless of whether they are accurate or not.
Users who cannot easily distinguish a correct answer from a confident but incorrect one
risk placing trust where it is not warranted; conversely, users who have experienced
model errors first-hand may develop excessive scepticism that prevents them from
benefiting from the technology's genuine capabilities.
The stakes are especially high in domains such as law and medicine, where decisions
grounded in incorrect information can carry serious personal consequences.
The opening question posed by legal practitioner Thomas Kuster at GenAI ZH 2026
captures precisely this tension: correctness alone is not sufficient; what matters is
whether a user can assess \emph{when} to trust a system's output.
 
Public awareness of this gap appears to be growing.
The same gfs.bern study found that the Swiss population identifies transparency,
accountability, and data security as the foremost conditions under which AI use
by public institutions would be considered acceptable.\cite{bern_chance_2025}
This demand for interpretability is not confined to the public sector.
Across domains, regulators, practitioners, and researchers are increasingly asking
how AI systems can be made to explain their reasoning in terms that non-expert users
can act on.
 
Explainable AI (XAI) has emerged as the principal research response to this challenge.
XAI encompasses methods and techniques designed to make AI behaviour more
interpretable, spanning post-hoc attribution methods such as SHAP and LIME,
intrinsic approaches such as Chain-of-Thought (CoT) prompting, and
human-centred techniques that leverage LLMs' own language generation capabilities
to produce natural language explanations \cite{wu_usable_2025}.
Yet despite the breadth of the XAI literature, a meaningful gap remains:
most existing work is oriented towards technical audiences, and there is limited
empirical evidence on which explanation types are most effective, or most helpful, for
lay users interacting with LLMs in everyday settings \cite{wu_usable_2025}.
 
\section{State of the Art}
 
Research into LLM explainability has expanded rapidly in response to the broad
deployment of these systems.
Existing approaches can be broadly categorized along three axes.
 
\textbf{Post-hoc attribution methods}, most notably SHAP \cite{lundberg_unified_2017} and LIME\cite{ribeiro_why_2016},
treat the model as a black box and assign importance scores to input features after
a prediction has been made.
Recent work has extended gradient-based, perturbation-based, surrogate-based, and
decomposition-based techniques to the token level, allowing individual words or phrases
to be highlighted as drivers of a model's output \cite{wu_usable_2025, smilkov_smoothgrad_2017, ancona_towards_2018, zhou_learning_2015, selvaraju_grad-cam_2020, binder_layer-wise_2016, grosse_studying_2023}.
 
\textbf{Intrinsic interpretability methods} embed transparency into the model's own generation process.
While rule-based models such as Generalized Additive Models \cite{caruana_intelligible_2015},
Explainable Boosting Machines \cite{lou_accurate_2013}, and Interpretable Decision Sets
\cite{lakkaraju_interpretable_2016} achieve transparency by design, these approaches are
fundamentally incompatible with the dense, high-dimensional parameter spaces of neural
networks; they cannot be applied to LLMs without abandoning the representational capacity
that makes such models effective.
Chain-of-Thought prompting is the most prominent example: by instructing a model
to articulate its reasoning step by step before arriving at a conclusion, CoT makes
the inferential process visible as part of the response itself \cite{wu_usable_2025}. 
Both ChatGPT \cite{guan_monitoring_2025} and Claude \cite{lindsey_biology_nodate} surface step-by-step reasoning traces as a standard transparency mechanism, making CoT the most widely deployed explanation paradigm
in production systems today.
More technically ambitious variants operate directly on the model's internal hidden-state
representations. For example, SelfIE \cite{chen_selfie_2024} enables LLMs to
interpret their own embeddings in natural language without additional training, but
such approaches require white-box access to model weights and remain opaque to
non-technical audiences.
 
\textbf{Human-centred and narrative explanations} represent a third strand, in which
LLMs generate natural language accounts of their own behaviour, either through
direct self-explanation or through a dedicated secondary model \cite{wu_usable_2025,you_gamifying_2024}.
A notable sub-class are counterfactual explanations, which characterize a model's
decision boundary by describing the minimal change to an input that would alter
the output \cite{wachter_counterfactual_2018}; extensions of this idea have
moved beyond single-feature perturbations toward actionable, multi-step
recourse paths that are more useful in practice \cite{rawal_beyond_2020}.

Gamification frameworks have further explored how narrative and interactive
presentation can increase engagement with explanations among lay users
\cite{you_gamifying_2024}.
 
\textbf{Human-AI Collaboration (HAIC)} research has drawn attention to additional dimensions
of the problem, emphasizing that effective collaboration depends not only on the quality
of explanations but also on mutual understanding, appropriate role allocation,
and calibrated trust between the human and the system
\cite{jeong_design_2025}.
Frameworks such as Human-in-the-Loop (HITL) and Human-on-the-Loop (HOTL)
position the user as an active participant in AI-assisted decision making, foregrounding
the need for tools that help users maintain appropriate oversight rather than simply
delegating judgment to the model.
 
A clear gap exists in the current literature.
While the technical inventory of XAI methods is rich, empirical evidence on which
approaches actually improve lay users' ability to calibrate their trust in LLM outputs
remains sparse.
The majority of evaluation studies focus on expert users, interpretability benchmarks,
or controlled classification tasks rather than on the open-ended, high-stakes question
formats that characterize everyday LLM use.
It is precisely this gap, namely the question of what XAI can do \emph{in practice} for
non-specialist users, that the present work addresses.
 
\section{Motivation}
 
The motivation for this work rests on three converging pressures.
 
First, \textbf{regulatory pressure} is mounting.
The EU AI Act, which entered into force in 2024, establishes explicit transparency and
explainability requirements for high-risk AI applications, and broader voluntary
frameworks from standardization bodies and industry coalitions are following suit.
Compliance requires not only that explanations exist, but that they are usable
by the people affected by AI decisions.
 
Second, \textbf{adoption is outpacing understanding}.
As documented above, more than half of Swiss residents who are aware of LLMs use
them in a professional context \cite{bern_chance_2025}, yet the majority of these
users have no formal training in machine learning.
The combination of high usage, high stakes, and low technical literacy creates
conditions in which miscalibrated trust (in either direction) can cause real harm.
 
Third, \textbf{trust is not self-correcting}.
Research on sycophancy in LLMs shows that models trained to be agreeable tend to
produce outputs that increase user trust by telling users what they want to hear,
independent of factual accuracy \cite{wu_usable_2025}.
Without external signals that make model uncertainty or evidential grounding visible,
users have little basis for adjusting their reliance on a given output.
XAI methods represent the most promising available mechanism for providing
such signals, and understanding which of them work for whom and in what
context is therefore a practically urgent question.
 
\section{Research Questions and Thesis Contribution}
\label{sec:research_questions_and_thesis}
 
Building on the motivation outlined above, this thesis addresses the following
research questions:
 
\begin{enumerate}
    \item Does the provision of explainability information help users to appropriately
          trust or distrust the output of an LLM?
    \item How do different XAI methods (specifically post-hoc attribution with SHAP and
          counterfactual reasoning, Chain-of-Thought prompting, and confidence
          indication) differentially affect lay users' trust levels?
    \item Is XAI a reliable indicator of whether the model is actually functioning
          correctly, or can it create misleading impressions of transparency?
    \item How can general users be made more aware of the appropriate degree of trust
          to place in AI-generated information?
    \item What is the computational and usability overhead of add-on XAI methods,
          and does this overhead limit their practical deployment?
\end{enumerate}
 
To answer these questions, the thesis makes the following contributions.
It reviews the relevant theoretical foundations spanning LLM architectures, XAI taxonomies, and the psychology of trust in AI systems.
It selects and implements three XAI approaches (SHAP-based token attribution combined
with counterfactual reasoning, Chain-of-Thought prompting, and confidence
indication) in a functional web-based demonstrator running on an open-source LLM.
It designs and conducts an empirical user study with \(n = 175\) prior LLM users,
measuring trust across five dimensions before and after exposure to each explanation
method.
Finally, it synthesizes the results into a prototype that combines the two
best-performing methods and derives concrete design recommendations for
trustworthy LLM-based applications.
 
\section{Thesis Structure}
 
The remainder of this thesis is organized as follows

\begin{itemize}
    \item \hyperref[ch:theory]{Chapter 2} provides the theoretical background, covering the architecture and life-cycle of LLMs, the taxonomy of XAI methods, and the multidimensional construct of trust in AI systems.
    \item \hyperref[ch:methods]{Chapter 3} describes the methodology, including the selection of the LLM and XAI methods, the implementation of the demonstrator, and the design of the user study.
    \item \hyperref[ch:results]{Chapter 4} presents the empirical results, including sample characteristics, trust outcomes per explanation method, and perceived helpfulness rankings.
    \item \hyperref[ch:prototype]{Chapter 5} Describes the implementation of a prototype that using a combination of the best rated XAI methods.
    \item \hyperref[ch:discussion]{Chapter 6} Discusses the findings, limitations, and future research directions.
\end{itemize}